\chapter{视觉模型}
\begin{quote}\small
\textit{本章摘要。}本章围绕图像中的缺陷识别展开。图像分类判断是否存在缺陷，目标检测给出缺陷的位置，图像分割描出缺陷的边界。先以卷积网络和 ResNet 解释如何从像素形成类别判断，再以 ViT 解释图像区域之间的注意力交互，并说明分类正确为何不等于局部特征对齐。随后以 Faster R-CNN 和 U-Net 分别讲解框级定位与像素级预测，并由它们的局限引出单阶段检测、多尺度特征和实例分割。缺陷样本稀少时，再用 DINOv2 系列的冻结图像块特征与少量正常样本逐块比较，做少样本异常检测。每个模型都沿输入、计算、输出和训练目标展开，最后回到细小缺陷、标注不足与跨相机验证这些典型问题。
\end{quote}
风险预警告诉维护人员何时需要关注设备，却不能指出缺陷位于哪里。本章转向巡检图像和表面质检。所用的图像级、框级或像素级标签都与未来故障标签不同。拆检后才能获得的照片可以支持离线诊断，却不能直接作为拆检前预警的输入。同样一组深浅像素，排列成细长的连续区域时可能表示裂纹，散布在表面时却可能只是纹理。模型既要识别这些空间关系，也要按任务要求保留位置。本章始终以表面质检为例，逐步回答三个问题。图中有没有裂纹，裂纹在哪里，哪些像素属于裂纹。

\section{视觉任务与标签粒度}
图像分类为整幅图给出类别，例如“正常”“裂纹”“划痕”，模型输出各类得分或概率。若裂纹与划痕可以同时出现，就需要为各类别分别标注是否存在，并使用多标签输出。分类适合工件筛查，但类别概率不能直接提供裂纹坐标。

目标检测为图中每个目标输出类别和矩形框，适合寻找多个缺陷、统计数量和指引复检位置。图像分割为每个像素给出标签，适合描述不规则边界和计算缺陷面积；语义分割只区分像素类别，实例分割还区分不同目标，同类的两条裂纹需要不同的实例编号。像素长度和面积要换算为物理尺寸，还需要相机标定或已知尺度。

三类任务分别需要图像级标签、框标注和像素标注，标签粒度决定了能直接训练什么输出：只有“有无裂纹”的标签时，不能直接计算逐像素分割损失。缺陷样本少到难以训练这些任务时，还可以只用少量正常样本建立参照，把偏离参照的位置报为异常，见第\ref{sec:vision-fewshot-ad}节。

这些任务通常共享一个提取图像特征的网络主体，称为\textbf{骨干网络}；骨干之后的\textbf{任务头}将特征转成类别、框或掩码。卷积神经网络（Convolutional Neural Network，CNN）与视觉 Transformer（Vision Transformer，ViT）是两种常用的视觉骨干。下面先通过分类走通从像素到预测的过程，再解释定位任务为何需要改变这个过程。
\begin{figure}[htbp]
\centering
\begin{tikzpicture}[node distance=0.65cm and 0.65cm]
 \node[idi input, minimum width=1.7cm] (image) {工业图像\\$H\times W$};
 \node[idi process, right=of image, minimum width=2.0cm] (backbone) {视觉骨干\\CNN / ViT};
 \node[idi state, right=of backbone, minimum width=2.0cm] (features) {空间特征\\单尺度或多尺度};
 \node[idi output, right=of features, minimum width=1.8cm] (class) {分类\\图像级};
 \node[idi output, below=0.75cm of class, minimum width=1.8cm] (detect) {检测\\框级};
 \node[idi output, above=0.75cm of class, minimum width=1.8cm] (segment) {分割\\像素级};
 \draw[idi arrow] (image)--(backbone); \draw[idi arrow] (backbone)--(features);
 \draw[idi arrow] (features)--(class); \draw[idi arrow] (features)--(detect); \draw[idi arrow] (features)--(segment);
\end{tikzpicture}
\caption{视觉系统先由卷积或 Transformer 骨干提取表示，再按标签粒度输出分类、检测或分割结果。}
\label{fig:vision-output-granularity}
\end{figure}
图\ref{fig:vision-output-granularity}中，分类可以汇总位置，检测与分割则要继续使用空间对应关系；原始 ViT 通常只输出单一尺度的图像块表示。

\section{图像分类与卷积骨干}
\subsection{卷积如何从像素中提取特征}
设彩色图像为$X\in\R^{H\times W\times C_{\rm in}}$，其中$H,W$为像素高度和宽度，$C_{\rm in}=3$时三个输入通道分别表示红、绿、蓝；隐藏层的通道表示学习到的特征，不再对应颜色。张量可以理解成带多个索引的数表，本章采用高、宽、通道的轴顺序。

若把图像直接展平后送入全连接层，每个位置都拥有独立权重，却没有利用“同一种局部形状可能出现在不同位置”这一规律。卷积用一个小窗口读取局部像素，并在整幅图上重复使用同一组权重。不同卷积核响应不同的局部变化，训练决定哪些响应对识别缺陷有用。

\paragraph{一个窗口怎样变成一个输出。}
取单通道窗口$\left[\begin{smallmatrix}1&2\\3&4\end{smallmatrix}\right]$与核$\left[\begin{smallmatrix}1&0\\0&-1\end{smallmatrix}\right]$，偏置为零，则对应输出为$1\times1+2\times0+3\times0-4\times1=-3$。核移到右边后仍使用这四个权重，只更换被读取的像素。若随后使用整流线性单元（Rectified Linear Unit，ReLU），即$\max(0,z)$，这一响应变为零。实际的核权重由训练学到，示例核只用于理解计算。

步幅是相邻窗口起点移动的像素数，填充是在图像边缘补值，空洞率是核内相邻采样点的间隔。例如$5\times5\times3$输入经过$3\times3$核、步幅一、无填充、四个输出通道，得到$3\times3\times4$，含$3\times3\times3\times4+4=112$个参数。每个输出通道都汇总三个输入通道，并不是分别对红绿蓝得到三个互不交流的答案。

\paragraph{卷积的索引、参数与等变性。}
设输入$X\in\R^{H\times W\times C_{\rm in}}$，方形核大小$k$，输出通道数$C_{\rm out}$，步幅$s$、填充$p$、空洞率$d_a$。采用从零开始的空间索引，有
\begin{equation}
 Z_{i,j,c}=b_c+\sum_{u,v=0}^{k-1}\sum_{a=1}^{C_{\rm in}}
 K_{u,v,a,c}X_{si-p+d_au,\,sj-p+d_av,\,a},
\end{equation}
越界输入按指定规则填充。这个深度学习中惯称的卷积实际采用互相关索引；若使用数学卷积，只需翻转核。沿高度方向，为使最下端采样不超过填充后的输入，需满足$s i+d_a(k-1)\leq H+2p-1$，因此
\begin{equation}
 H_{\rm out}=\left\lfloor\frac{H+2p-d_a(k-1)-1}{s}\right\rfloor+1,
\end{equation}
宽度同理。参数量为$k^2C_{\rm in}C_{\rm out}+C_{\rm out}$，与图像面积无关；不共享的局部连接则须为每个输出位置另存一组权重。

在无限网格或忽略边界、$s=1$时，令$(T_\delta X)_{i,j}=X_{i-\delta_1,j-\delta_2}$，直接代入求和式得$K*(T_\delta X)=T_\delta(K*X)$。这就是平移等变性，而不是平移后输出保持不变。步幅$s>1$时通常仅保留对$s$整数倍平移的对应关系，边界填充还会引入额外差异。图像分类可以通过汇聚追求不变性，定位任务却应保留位置随输入一起移动的关系。

\subsection{多层卷积如何形成整幅图的判断}
单个小窗口只能看到裂纹的一部分。把卷积与非线性激活连续堆叠后，后层可以组合前层的局部响应，例如把多个方向一致的边缘组合为较长的结构。网络常用步幅大于一的卷积或池化降低空间分辨率，这个过程称为下采样；它减少计算量，也可能让细裂纹与背景混在一起。

一个深层位置在输入中能够访问的范围称为\textbf{理论感受野}。令$r_l$表示第$l$层的感受野宽度，$j_l$表示该层相邻位置在原图中的间隔。以$r_0=j_0=1$开始，对于核大小$k_l$、步幅$s_l$和空洞率$d_l$的层，有
\begin{equation}
 r_l=r_{l-1}+d_l(k_l-1)j_{l-1},\qquad j_l=s_lj_{l-1}.
 \label{eq:vision-receptive-field}
\end{equation}
新增的有效间隔为$d_l(k_l-1)$个输入位置，每个位置跨越$j_{l-1}$个原图像素。三层$3\times3$卷积的空洞率均为$1$、步幅依次为$1,2,1$时，$(r_l,j_l)$依次为$(3,1),(5,2),(9,2)$。第三层每个内点理论上访问九像素宽的区域，相邻输出却只相隔两像素。这个范围不表示每个像素贡献相同，边界处还可能读到填充值。空洞率大于$1$时，访问范围扩大，但卷积核仍只采样其中的离散位置。

设最后一层特征为$F\in\R^{H'\times W'\times D}$，其中$D$为通道数。分类可以先对每个通道作全局平均，再进行线性分类：
\begin{equation}
 v_c=\frac{1}{H'W'}\sum_{i=1}^{H'}\sum_{j=1}^{W'}F_{i,j,c},\qquad
 z=Uv+b,\qquad p_k=\frac{e^{z_k}}{\sum_{a=1}^{K}e^{z_a}}.
 \label{eq:vision-classification}
\end{equation}
这里$K$为互斥类别数，$v\in\R^D$为整图表示，$U\in\R^{K\times D}$与$b\in\R^K$为分类头参数，$z_k$是类别$k$的实数得分，$p_k$是模型给出的类别概率。softmax 假设每幅图只有一个真实类别。若真实类别为$y$，单幅图的交叉熵损失为$-\log p_y=-z_y+\log\sum_{a=1}^K e^{z_a}$，对各得分求导得
\[
 \frac{\partial(-\log p_y)}{\partial z_k}=p_k-\mathbf1_{k=y}.
\]
例如三类概率为$(0.7,0.2,0.1)$而真实类别为第二类时，梯度为$(0.7,-0.8,0.1)$：沿负梯度方向更新会增大真实类别的得分，并按各自概率减小其余类别的得分。这一梯度再经反向传播调整分类头和卷积核，使特征对区分类别有用。

全局平均把各位置汇总为一个向量，适合回答“有没有裂纹”，却舍弃了直接读取位置的接口；检测与分割必须在汇总之前保留空间特征。

\subsection{网络加深以后的训练}
增加层数能让模型组合更大范围的局部信息，但普通深层网络可能比浅层网络具有更高的训练误差。连训练集都拟合得更差，就不能只用过拟合解释：更长的连续变换使优化更困难，深层网络也不容易自行学出“保持已有表示不变”的操作。

残差网络（Residual Network，ResNet）将一个块的输出写为\footnote{来源：He 等，\href{https://arxiv.org/abs/1512.03385}{\textit{Deep Residual Learning for Image Recognition}}，2016。}
\begin{equation}
 h_{l+1}=h_l+F_l(h_l).
 \label{eq:vision-residual}
\end{equation}
$h_l$是输入特征，$F_l$是由卷积、归一化和激活组成的残差分支。这里先省略相加后的激活，以分析加法节点的梯度；原始 ResNet 在相加后还施加 ReLU，完整块的导数还须乘入该激活的导数。若期望映射为$H(h_l)$，残差分支要学习的就是差值$H(h_l)-h_l$；已有表示已经合用时，分支输出接近零即可在加法节点保留输入。图\ref{fig:vision-residual-block}展示了分支与捷径如何在同一位置相加。

\begin{figure}[htbp]
\centering
\begin{tikzpicture}[node distance=0.8cm]
 \node[idi input] (h) {$h_l$};
 \node[idi process,right=of h,minimum width=3cm] (f) {卷积、归一化、激活\\$F_l(h_l)$};
 \node[idi state,right=of f] (add) {$+$};
 \node[idi output,right=of add] (out) {$h_{l+1}$};
 \draw[idi arrow] (h)--(f); \draw[idi arrow] (f)--(add); \draw[idi arrow] (add)--(out);
 \draw[idi arrow] (h.north)--++(0,0.8) -| node[pos=0.25,above,font=\small]{恒等捷径} (add.north);
\end{tikzpicture}
\caption{同尺寸残差块的计算路径。捷径传递输入，分支学习修正量，二者逐元素相加。图中省略了具体实现可能在相加后使用的激活。}
\label{fig:vision-residual-block}
\end{figure}

\paragraph{残差梯度不再只有单一路径。}
设各层$h_l\in\R^d$，且残差分支$F_l:\R^d\to\R^d$可微。令$J_l=\partial F_l/\partial h_l^{\mathsf T}$，则
\begin{align}
 \frac{\partial h_{l+1}}{\partial h_l^{\mathsf T}}&=I+J_l,\\
 \nabla_{h_l}\mathcal L&=(I+J_l)^{\mathsf T}\nabla_{h_{l+1}}\mathcal L.
\end{align}
连续两个块的前向Jacobian为$(I+J_{l+1})(I+J_l)=I+J_l+J_{l+1}+J_{l+1}J_l$，包含不经过任何残差分支的恒等项、经过一个分支的项和经过两个分支的项。相比普通网络只有$J_{l+1}J_l$，残差结构允许梯度走较短路径。以标量为例：设每个分支的导数都是$0.1$，普通网络把$10$层的导数连乘得$0.1^{10}=10^{-10}$，残差网络每层导数为$1+0.1$，连乘约为$2.6$，梯度量级不随层数增加而迅速衰减。若$\|J_l\|_2\leq\epsilon<1$，则对任意向量$v$有
\begin{equation}
 (1-\epsilon)\|v\|_2\leq\|(I+J_l)v\|_2\leq(1+\epsilon)\|v\|_2.
\end{equation}
小残差分支使单块接近恒等，但多块的上下界仍相乘，所以残差不是绝不消失或爆炸的定理。改变通道或空间尺寸时，要把捷径替换为投影$S_lh_l$，导数变为$S_l+J_l$，此时也不再有字面上的恒等路径。

\paragraph{相加要求尺寸一致。}
若残差输入和分支输出都是$8\times8\times16$，逐元素相加后仍是$8\times8\times16$。若残差分支改变空间尺寸或通道数，捷径也要用步幅卷积等投影变换得到相同尺寸。后文 U-Net 将使用通道拼接，同样两组特征拼接后为$8\times8\times32$，两种连接的输出维度和计算含义不同。

\subsection{从基本块到瓶颈块与轻量卷积}
ResNet-18 和 ResNet-34 主要使用由两个$3\times3$卷积构成的基本残差块。ResNet-50 等更深的版本使用瓶颈块：先以$1\times1$卷积减少通道，在较窄的表示上做$3\times3$卷积，再以$1\times1$卷积恢复通道。$1\times1$卷积只混合同一位置的不同通道，瓶颈设计把较昂贵的空间卷积放在较少通道上，使加深网络的计算代价更可控。

图\ref{fig:vision-resnet-original-blocks}左侧基本块维持 64 通道，右侧瓶颈块先将 256 通道压到 64，再恢复为 256。方框内依次标出核大小与输出通道数，弧线是恒等捷径。两种块均在相加后施加 ReLU，属于原始版本；该图没有单列卷积配套的批归一化。

\begin{figure}[H]
\centering
\includegraphics[width=0.84\linewidth]{images/中文重绘/resnet-cn.png}
\caption[ResNet 残差块]{ResNet 残差块\protect\footnotemark}
\label{fig:vision-resnet-original-blocks}
\end{figure}
\footnotetext{来源：He 等，\href{https://arxiv.org/abs/1512.03385v1}{\textit{Deep Residual Learning for Image Recognition}}，2016。}

若相机侧计算资源更紧，MobileNet 一类轻量网络使用深度可分离卷积：先逐通道做空间卷积，再用$1\times1$卷积混合通道。忽略偏置，标准卷积每个位置约需$k^2C_{\rm in}C_{\rm out}$次乘加，拆分后为$k^2C_{\rm in}+C_{\rm in}C_{\rm out}$，二者比值为$1/C_{\rm out}+1/k^2$；$k=3$、$C_{\rm out}=64$时约为$0.13$，乘加数降到约八分之一。节省来自对空间与通道混合方式的限制，表达能力也随之改变，实际延迟还要在目标设备上测量。

这些设计都没有取消局部卷积的约束。若判断一个局部缺陷必须参考远处的部件，能否让区域之间直接交换信息？

\section{ViT 与区域注意力}
\subsection{从图像块到类别预测}
卷积通过逐层组合扩大观察范围。视觉 Transformer 将图像划分为一组图像块（patch），把每个块编码为一个向量单元（token，下文称标记），再通过自注意力更新这些向量\footnote{来源：Dosovitskiy 等，\href{https://arxiv.org/abs/2010.11929}{\textit{An Image is Worth 16x16 Words: Transformers for Image Recognition at Scale}}，2021。}。它仍可用于整图分类，改变的是骨干内部组织空间信息的方法。

设输入为$H\times W\times C$，每块大小为$P\times P$，且$H,W$均可被$P$整除。共有$N=HW/P^2$个块。第$i$块展平为列向量$x_i\in\R^{P^2C}$，线性投影为
\begin{equation}
 z_i=Ex_i+e_i^{\rm pos},\qquad E\in\R^{d\times P^2C}.
 \label{eq:vision-patch-embedding}
\end{equation}
$z_i\in\R^d$为块表示，$e_i^{\rm pos}\in\R^d$为位置嵌入。投影权重在各块之间共享，位置嵌入标明块来自哪里；没有位置信息时，普通自注意力本身不区分相同图像块在网格中的排列。

一幅$32\times32\times3$图像按$8\times8$划块，得到$16$个块，每块包含$192$个数。投影到$d=64$后，按行放置这些向量得到$Z\in\R^{16\times64}$。对单头注意力，令$W_Q,W_K,W_V\in\R^{64\times64}$，计算
\begin{equation}
 Q=ZW_Q,\quad K=ZW_K,\quad V=ZW_V,\qquad
 A=\softmax\left(\frac{QK^{\mathsf T}}{\sqrt{64}}\right),\quad O=AV.
 \label{eq:vision-attention}
\end{equation}
$Q,K,V$分别称为查询、键和值，$A\in\R^{16\times16}$沿每一行归一化。$A_{ij}$决定更新第$i$块时读取第$j$块值向量的权重，输出$O$仍为$16\times64$。每个块因此可以在同一层读取任意其他块的信息；权重由当前表示计算，不等于人工给定的空间邻接关系。

原始 ViT 在块序列前加入一个可学习的分类标记，此时本例共有$17$个位置。经过多层带残差连接的注意力与逐位置前馈网络后，取分类标记的最终表示送入分类头，用图像标签计算交叉熵；也可以改用块表示的平均值作分类。用于检测或分割时，则要保留图像块与原网格的对应关系，另外设计空间输出头。

图\ref{fig:vision-vit-original}中$0*$是额外学习的分类标记，不由图像块投影得到；其余标记由图像块投影并加入位置嵌入形成。右侧是重复堆叠的前置归一化编码块，分类头读取分类标记的最终表示。图中只画出九个图像块，与前述十六块算例的数量不同，计算规则一致。

\begin{figure}[H]
\centering
\includegraphics[width=0.94\linewidth]{images/vit-cn-v2.png}
\caption[ViT 分类结构]{ViT 分类结构\protect\footnotemark}
\label{fig:vision-vit-original}
\end{figure}
\footnotetext{来源：Dosovitskiy 等，\href{https://arxiv.org/abs/2010.11929}{\textit{An Image is Worth 16x16 Words: Transformers for Image Recognition at Scale}}，2021。}

\FloatBarrier
\subsection{空间特征伪影与寄存标记}
\label{sec:vision-register-artifacts}
图像级分类监督只要求分类标记给出整图类别，并不直接约束每个图像块都保留准确的局部语义。在一些视觉 Transformer 中，低信息背景区域会出现范数很大的异常标记，在中间层承担存储或传递全局信息的作用。Registers 论文在输入序列中加入额外的可学习标记，给这种内部计算提供专门位置，输出时丢弃它们，只保留分类标记和图像块表示\footnote{来源：Darcet 等，\href{https://arxiv.org/abs/2309.16588}{\textit{Vision Transformers Need Registers}}，2024。}。可将输入写成
\begin{equation}
 Z_{\rm in}=[z_{\rm CLS},z_1,\ldots,z_N,r_1,\ldots,r_R]^{\mathsf T}
 \in\R^{(1+N+R)\times d},
 \label{eq:vision-register-input}
\end{equation}
其中每一行对应一个标记，$R$为寄存标记数。寄存标记不对应原图上的固定位置，在训练中与模型共同学习。图\ref{fig:vision-registers}中蓝色为图像块标记、黑色为分类标记、黄色为寄存标记，顶部两个括号分别标出保留和丢弃的输出；图中寄存标记的编号沿用原论文写到$N$，对应式\eqref{eq:vision-register-input}中的$R$个寄存标记。

后续研究认为，减少高范数标记并不必然消除空间语义错位：全局注意力和粗粒度的图像级监督可能使背景图像块成为表达全局语义的捷径，导致分类正确而局部特征与目标位置对齐较差\footnote{来源：Shi 等，\href{https://arxiv.org/abs/2602.22394}{\textit{Vision Transformers Need More Than Registers}}，2026。}。工业质检因此要用缺陷框、掩码等定位指标检查，注意力图只能提供诊断线索。

\begin{figure}[H]
\centering
\includegraphics[width=0.98\linewidth]{images/registers-cn-v2.png}
\caption[ViT 的寄存标记]{ViT 的寄存标记\protect\footnotemark}
\label{fig:vision-registers}
\end{figure}
\footnotetext{来源：Darcet 等，\href{https://arxiv.org/abs/2309.16588}{\textit{Vision Transformers Need Registers}}，2024。}

\FloatBarrier
\subsection{为什么引入窗口注意力与层级表示}
\label{sec:vision-window}
全局注意力显式形成$N\times N$的得分矩阵，主要交互计算量为$O(N^2d)$。若高、宽都加倍而块大小不变，块数增加到四倍，注意力得分数量增加到十六倍。细缺陷需要较高分辨率，全局交互的成本因此成为直接限制。

Swin Transformer 将注意力限制在局部窗口内，并在相邻层之间移动窗口划分：前一层位于窗口边缘的两个块，下一层可以进入同一窗口，跨窗口信息由此逐层传递。若每个窗口含$M\times M$个单元且可整分，则有$N/M^2$个窗口，交互计算量为
\begin{equation}
 O\left(\frac{N}{M^2}(M^2)^2d\right)=O(NM^2d).
\end{equation}
固定$M$时，这部分成本随$N$线性增加，代价是单层只能读取窗口内信息。Swin 还逐级合并相邻单元，得到不同空间尺度的特征，供检测与分割使用；合并降低成本，也可能损失局部细节。

图\ref{fig:vision-swin-original-windows}用红框标出局部注意力窗口，灰色小格表示图像块。第$l+1$层改变窗口边界，使原先分属两个窗口的块能够在新窗口内交互；移动的是分组边界，不是把图像内容平移。

\begin{figure}[H]
\centering
\includegraphics[width=0.88\linewidth]{images/中文重绘/swin-cn.png}
\caption[Swin 窗口注意力]{Swin 窗口注意力\protect\footnotemark}
\label{fig:vision-swin-original-windows}
\end{figure}
\footnotetext{来源：Liu 等，\href{https://arxiv.org/abs/2103.14030v2}{\textit{Swin Transformer: Hierarchical Vision Transformer using Shifted Windows}}，2021。}

卷积把局部性和权重共享直接写入结构。原始 ViT 的这类约束较弱，通常更依赖预训练规模和训练策略。实际比较应控制输入分辨率、预训练数据和微调预算。无论哪种骨干，输入缩放或块投影已经丢失的信息都不能被自动补回。块投影在$d<P^2C$时不可能保留所有输入变化，但这不意味着小于块的缺陷一定不可见，结果取决于投影保留了什么以及标签要求模型识别什么。

\section{目标检测与框级定位}
\subsection{为什么先提议区域，再分类和修正位置}
分类器每图输出一组类别概率，检测器却需要输出数量不固定的一组目标：一张图可能没有缺陷，也可能有多条裂纹。若把分类器在所有可能的矩形上重复运行，不仅计算量大，相邻矩形还会反复读取相同像素。

Faster R-CNN 沿用区域卷积神经网络（Region-based Convolutional Neural Network，R-CNN）的区域识别思路，把主要计算集中在共享骨干中，分两阶段工作：区域提议网络（Region Proposal Network，RPN）在共享特征上为预置参考框预测“目标还是背景”以及坐标偏移，筛选出候选区域；第二阶段对候选区域提取固定大小的特征，输出具体类别和进一步修正的框。候选区域也可能只包含背景，因此第二阶段仍需背景类。

\begin{figure}[htbp]
\centering
\begin{tikzpicture}[node distance=0.6cm]
 \node[idi input] (image) {图像};
 \node[idi state,right=of image] (feat) {共享特征};
 \node[idi process,right=of feat] (rpn) {RPN\\候选框};
 \node[idi process,right=of rpn] (roi) {区域特征\\固定大小};
 \node[idi output,right=of roi] (head) {类别\\修正框};
 \draw[idi arrow] (image)--(feat); \draw[idi arrow] (feat)--(rpn);
 \draw[idi arrow] (rpn)--(roi); \draw[idi arrow] (roi)--(head);
 \draw[idi arrow] (feat.south)--++(0,-0.7) -| node[pos=0.25,below,font=\small]{按候选位置读取特征} (roi.south);
\end{tikzpicture}
\caption{Faster R-CNN 的两阶段计算。区域提议给出在哪里读取，共享特征提供读取内容，第二阶段进行区域分类与框修正。}
\label{fig:vision-faster-rcnn}
\end{figure}

图\ref{fig:vision-faster-rcnn}中的区域特征来自整图共享特征，不必为每个候选重新运行骨干。原始 Faster R-CNN 用感兴趣区域（Region of Interest，RoI）池化把大小不同的区域转成固定大小的特征，其坐标量化会影响精细对齐，后文 Mask R-CNN 的 RoIAlign 用插值采样减轻这一问题。

\subsection{框怎样匹配标签并参与训练}
锚框（anchor）是在特征位置预置的参考矩形，不是真实标注。若一幅图有$100$个候选、两类缺陷，含背景类的简化检测头输出$100\times3$的类别得分和$100\times4$的类别无关框偏移。训练先把候选与标注匹配，得到“这是背景还是哪类缺陷”及正样本的坐标目标，再计算损失；不同检测器的输出格式可能不同。

图\ref{fig:vision-rpn-original}展开 RPN 在一个特征位置的计算：滑动窗口形成中间表示，两个分支分别输出$k$个锚框的$2k$个目标／背景得分与$4k$个框回归量；右侧矩形是不同尺度和宽高比的锚框，并非筛选后的检测结果。图中的 256 维对应原论文 ZF 骨干，采用 VGG 时中间表示为 512 维。

\begin{figure}[H]
\centering
\includegraphics[width=0.78\linewidth]{images/中文重绘/faster-rcnn-cn.png}
\caption[RPN 区域提议]{RPN 区域提议\protect\footnotemark}
\label{fig:vision-rpn-original}
\end{figure}
\footnotetext{来源：Ren 等，\href{https://arxiv.org/abs/1506.01497v1}{\textit{Faster R-CNN: Towards Real-Time Object Detection with Region Proposal Networks}}，2015。}

参考框或候选框与真实标注匹配时，常用的几何依据是交并比（Intersection over Union，IoU）。两个框$A,B$的交并比为
\begin{equation}
 \operatorname{IoU}(A,B)=\frac{|A\cap B|}{|A\cup B|}.
\end{equation}
这里$|\cdot|$表示面积，IoU 取值为$[0,1]$。匹配到目标的候选是正样本，背景候选是负样本，阈值之间的候选还可以按协议忽略。RPN 的类别目标是目标或背景，第二阶段则区分具体缺陷类别。下面用第二阶段的简化目标解释分类与定位如何一起学习，完整训练还包括 RPN 的提议损失。

\paragraph{从几何框到可训练的回归目标。}
设参考框中心及宽高为$(a_x,a_y,a_w,a_h)$，真实框为$(g_x,g_y,g_w,g_h)$，宽高均严格为正且使用同一像素坐标系。相同相对位置误差在大框中对应更大的像素位移。为减小回归目标对框尺寸的依赖，定义无量纲目标
\begin{equation}
 t_x^*=\frac{g_x-a_x}{a_w},\quad t_y^*=\frac{g_y-a_y}{a_h},\quad
 t_w^*=\log\frac{g_w}{a_w},\quad t_h^*=\log\frac{g_h}{a_h}.
\end{equation}
网络输出$t\in\R^4$后，以$b_x=a_x+a_wt_x,b_y=a_y+a_ht_y,b_w=a_we^{t_w},b_h=a_he^{t_h}$解码。中心相对化降低尺度依赖，对数宽高保证解码尺寸为正；但$\partial b_w/\partial t_w=b_w$，极端输出仍可能带来不稳定。

对匹配到真实框的正样本，常用残差$e=t-t^*$的分段损失
\begin{equation}
 \rho_\beta(e)=\begin{cases}e^2/(2\beta),&|e|<\beta,\\|e|-\beta/2,&|e|\geq\beta,\end{cases}
 \qquad \beta>0.
\end{equation}
在接点$e=\pm\beta$，两段函数值均为$\beta/2$，两侧导数也分别相等，所以此处分段不妨碍一阶求导。其导数在小误差处为$e/\beta$，大误差处为$\operatorname{sign}(e)$，兼顾局部平滑与异常残差的有界影响。完整损失可写为
\begin{equation}
 \mathcal L=\frac1{N_{\rm cls}}\sum_i\ell_{\rm cls}(p_i,y_i)
 +\frac{\lambda}{\max(1,N_+)}\sum_{i:y_i>0}\sum_{j\in\{x,y,w,h\}}\rho_\beta(t_{ij}-t_{ij}^*).
\end{equation}
这里$p_i$是候选$i$的类别概率向量，$y_i=0$表示背景、$y_i>0$表示缺陷类别；$N_{\rm cls}>0$为参与分类计分的候选数，$N_+$为正样本数，$\lambda\geq0$控制定位项权重。背景样本没有真实框回归目标；分类与定位的归一化方式必须报告，否则相同$\lambda$可能代表不同梯度比例。

IoU直接比较几何重叠。设两个边长为$a$的正方形仅水平错位$\delta\in[0,a)$，则交集为$a(a-\delta)$，并集为$a(a+\delta)$，所以$\operatorname{IoU}=(a-\delta)/(a+\delta)$，在$\delta=0$右侧的导数为$-2/a$。这说明同样一个像素的偏移对小框更加苛刻。不相交框的普通IoU为零，在保持不相交的区域内不给出靠近方向，因此实际定位还会结合参数回归或加入外接框、中心距离等项。

\subsection{推理怎样得到数量不固定的检测结果}
推理时没有真实框可以匹配。模型先根据参考框解码预测偏移，再按类别分数筛掉低分框。多个候选可能指向同一条裂纹，因此还需要消除重复结果。非极大值抑制（Non-Maximum Suppression，NMS）通常逐类进行：先保留最高分框，再删除与它的 IoU 超过阈值的其他框，对剩余候选重复这一过程。例如三个裂纹框的得分为$0.9,0.8,0.6$，后两个框与第一个框的 IoU 分别为$0.7$和$0.1$；阈值取$0.5$时，保留$0.9$的框、删去$0.8$的框，$0.6$的框与它几乎不重叠，得以保留。最终框数由图像和阈值共同决定；两个紧邻的真实缺陷也可能互相抑制。

检测评价同样需要空间匹配。先按置信度排序，同类预测框在给定 IoU 阈值下与真实框一对一匹配。未匹配预测计为误报，未匹配真实框计为漏检。平均精度（Average Precision，AP）汇总精确率与召回率的关系。平均精度均值（mean Average Precision，mAP）再按指定类别或 IoU 阈值取平均。不同协议的 mAP 不能直接比较。小框对偏移敏感，应另外报告不同缺陷尺寸下的召回。

\subsection{单阶段检测与候选识别}
YOLO（You Only Look Once）代表的单阶段检测直接从整图特征的多个位置预测类别与坐标，不再为每个候选运行第二阶段识别头，适合吞吐量要求高的任务。早期 YOLO 由网格单元预测框、目标置信度和类别，密集目标和尺度变化较难处理；后续版本改进了多尺度输出与预测头，有的使用锚框，有的不用，使用时应按具体版本核对输出和正样本分配。

单阶段检测会面对大量背景位置。许多背景已经容易识别，却仍可能在总损失中占据较大比例。RetinaNet 用焦点损失（focal loss）降低易样本权重：
\begin{equation}
 \mathcal L_{\rm focal}=-\alpha_t(1-p_t)^\gamma\log p_t.
\end{equation}
二分类时，$p_t$是对真实标签的预测概率，即正样本取$p$，负样本取$1-p$。$\alpha_t$为类别权重，$\gamma\geq0$控制降权强度。例如$\gamma=2$时，$p_t=0.99$的易样本权重为$0.0001$，$p_t=0.5$的样本权重为$0.25$。这改变了不同难度样本对训练的相对影响，标注错误也可能因此获得较大权重。

\paragraph{焦点损失怎样改变梯度。}
对正类样本设logit为$z$、$p=\sigma(z)$，则$\ell=-\alpha(1-p)^\gamma\log p$。不能只把交叉熵梯度乘以$(1-p)^\gamma$，因为这个权重本身依赖参数。乘积法则给出
\begin{equation}
 \frac{d\ell}{dz}=\alpha\gamma p(1-p)^\gamma\log p-\alpha(1-p)^{\gamma+1}.
\end{equation}
令$p=1-\varepsilon$并利用$\log(1-\varepsilon)\approx-\varepsilon$，得梯度约为$-\alpha(\gamma+1)\varepsilon^{\gamma+1}$；普通交叉熵则约为$-\alpha\varepsilon$。所以易样本梯度以更高阶衰减，困难样本占据更大的相对权重。但标注错误往往也是困难样本，焦点损失不具备自动甄别错误标签的能力，也不能保证输出概率仍然校准。

\subsection{FPN 与多尺度特征}
小裂纹需要较密的空间采样，大范围表面损伤则需要较广的上下文；只在最深层特征上检测，可能已经失去小缺陷的位置。特征金字塔网络（Feature Pyramid Network，FPN）在骨干不同层的特征之间建立自顶向下的信息路径。

设骨干特征$C_l$的分辨率随层级$l$增加而降低。先用$1\times1$卷积$L_l$统一通道数，再从最高层开始融合：
\begin{equation}
 \widetilde P_L=L_L(C_L),\qquad
 \widetilde P_l=L_l(C_l)+\operatorname{Up}(\widetilde P_{l+1}).
 \label{eq:vision-fpn}
\end{equation}
$\operatorname{Up}$将高层特征上采样到$C_l$的空间大小，两个分支通道相同后逐元素相加，通常再用$3\times3$卷积处理融合结果，得到检测使用的$P_l$。高层提供较广的语义信息，横向连接提供较密的空间信息，不同大小的目标可分配到不同层级上预测。FPN 可以接入两阶段或单阶段检测器，它改变特征来源，不规定最终输出格式。

图\ref{fig:vision-fpn-original}中，左侧是骨干的自底向上路径，右侧是产生多尺度预测特征的自顶向下路径。下方放大框对应式\eqref{eq:vision-fpn}中的一次融合：高层特征上采样两倍，横向特征经过$1\times1$卷积，二者相加；融合后的$3\times3$卷积没有单列。

\begin{figure}[H]
\centering
\includegraphics[width=0.62\linewidth,height=0.44\textheight,keepaspectratio]{images/中文重绘/fpn-cn.png}
\caption[FPN 多尺度特征融合]{FPN 多尺度特征融合\protect\footnotemark}
\label{fig:vision-fpn-original}
\end{figure}
\footnotetext{来源：Lin 等，\href{https://arxiv.org/abs/1612.03144v2}{\textit{Feature Pyramid Networks for Object Detection}}，2017。}

上采样只是把数表放大。FPN 能利用的细节来自仍保留这些信息的横向分支；若缺陷在所有分支中都已丢失，融合不能把它恢复。这个限制同样存在于像素级分割中。

\section{图像分割与像素级预测}
\subsection{为什么解码器需要编码器的细节}
矩形框可以圈出裂纹，却同时包含很多正常背景，不能直接用于测量裂纹面积。分割将预测目标改为与图像对应的像素标签；对细长裂纹，模型既要判断局部变化是否属于裂纹，又要保持其边界与连续性。

U-Net 是具有编码器与解码器的卷积网络\footnote{来源：Ronneberger 等，\href{https://arxiv.org/abs/1505.04597}{\textit{U-Net: Convolutional Networks for Biomedical Image Segmentation}}，2015。}。编码器逐级下采样，形成越来越粗的空间特征；解码器逐级上采样，将表示恢复到输出网格。只把最深层特征放大时，精细位置难以确定，因此 U-Net 还把编码器同尺度特征传给解码器。图\ref{fig:vision-unet}展示了这一计算关系。

\begin{figure}[htbp]
\centering
\begin{tikzpicture}[node distance=0.65cm and 2.0cm]
 \node[idi state] (e1) {编码特征\\$32\times32\times16$};
 \node[idi state,below=of e1] (e2) {编码特征\\$16\times16\times32$};
 \node[idi process,below=of e2] (b) {深层特征\\$8\times8\times64$};
 \node[idi process,right=of e2] (d2) {上采样、拼接、卷积\\$16\times16\times32$};
 \node[idi process,right=of e1] (d1) {上采样、拼接、卷积\\$32\times32\times16$};
 \node[idi output,above=of d1] (out) {像素概率\\$32\times32$};
 \draw[idi arrow] (e1)--(e2); \draw[idi arrow] (e2)--(b);
 \draw[idi arrow] (b.east)-|(d2.south); \draw[idi arrow] (d2)--(d1); \draw[idi arrow] (d1)--(out);
 \draw[idi arrow,dashed] (e2)--node[above,font=\small]{同尺度拼接}(d2);
 \draw[idi arrow,dashed] (e1)--node[above,font=\small]{同尺度拼接}(d1);
\end{tikzpicture}
\caption{采用尺寸保持卷积的简化 U-Net。实线表示下采样或解码计算，虚线把编码特征送到同尺度解码层。尺寸为教学设置，未逐层复现原论文配置。}
\label{fig:vision-unet}
\end{figure}

例如，把$8\times8\times64$深层特征上采样并调整为$16\times16\times32$，与编码器的$16\times16\times32$特征按通道拼接，得到$16\times16\times64$，再通过卷积产生$16\times16\times32$解码特征。拼接保留两组特征供后续学习使用，和 ResNet 的逐元素相加不同。图中为了保持尺寸使用了填充，原始 U-Net 的无填充卷积需要裁切编码特征后再拼接。

图\ref{fig:vision-unet-original}给出原始 U-Net 的完整尺寸关系：蓝色箭头表示$3\times3$无填充卷积与 ReLU，每次使高、宽各减少二；红色箭头下采样，绿色箭头上卷积，灰色箭头复制并裁剪编码特征。输入为$572\times572$，输出为$388\times388$的两类分割得分图，尚不是阈值化后的二值掩码。

\begin{landscape}
\special{pdf: put @thispage << /Rotate 90 >>}
\thispagestyle{plain}
\setlength{\intextsep}{2pt}
\captionsetup{skip=5pt}
\begin{figure}[H]
\centering
\includegraphics[width=0.95\linewidth,height=0.87\textwidth,keepaspectratio]{images/中文重绘/unet-cn.png}
\caption[U-Net 编码器与解码器]{U-Net 编码器与解码器\protect\footnotemark}
\label{fig:vision-unet-original}
\end{figure}
\footnotetext{来源：Ronneberger 等，\href{https://arxiv.org/abs/1505.04597v1}{\textit{U-Net: Convolutional Networks for Biomedical Image Segmentation}}，2015。}
\end{landscape}

同尺度连接使解码器可以同时读取细节与上下文，但也可能带入背景纹理；它提供了使用信息的路径，是否正确使用仍由训练决定。

\subsection{从概率图到掩码与训练损失}
以一张$32\times32$图像的二值分割为例，模型输出$32\times32$个logit，即尚未转成概率的实数得分，经sigmoid成为同形状概率图，训练时与同形状的零一标签比较，推理时才按验证集确定的阈值变成硬掩码。多类互斥分割通常输出$32\times32\times K$，在每个像素的$K$个类别上归一化，而不是在整幅图的全部像素上做一次softmax。

四个像素的真实掩码为$(1,1,0,0)$，预测硬掩码为$(1,0,1,0)$，则交集为一、并集为三，IoU为$1/3$，Dice为$2/(2+2)=1/2$。训练中的软Dice用概率代替硬预测，以便求导；阈值化通常不参与这条训练梯度。

\paragraph{像素似然、Dice与稀疏前景。}
设一幅图含$n$个像素，二值标签$y_i\in\{0,1\}$，模型输出$p_i=\sigma(z_i)$，其中$\sigma(z)=1/(1+e^{-z})$为 sigmoid 函数。将每个像素的标签建模为伯努利分布，在给定图像后采用乘积形式得到似然。取负对数并按像素数平均，得到二值交叉熵（Binary Cross Entropy，BCE）：
\begin{align}
 p(y\mid X)&=\prod_{i=1}^n p_i^{y_i}(1-p_i)^{1-y_i},\\
 \mathcal L_{\rm BCE}&=-\frac1n\sum_i[y_i\log p_i+(1-y_i)\log(1-p_i)],\\
 \frac{\partial\mathcal L_{\rm BCE}}{\partial z_i}&=\frac{p_i-y_i}{n}.
\end{align}
这个似然分解不表示真实裂纹像素相互独立，而是输出头采用的近似。前景极少时，大量背景项会主导平均损失。Dice从集合重叠出发：硬掩码的$D=2|A\cap B|/(|A|+|B|)$，将预测指示函数换成概率便得到软 Dice 损失：
\begin{equation}
 \mathcal L_{\rm Dice}=1-\frac{2\sum_i p_iy_i+\epsilon}{\sum_i p_i+\sum_i y_i+\epsilon}.
 \label{eq:vision-soft-dice}
\end{equation}
其中$\epsilon>0$是固定的小平滑常数，用于避免空掩码时分母为零，求导时不随$p_i$变化。记$N=2\sum_i p_iy_i+\epsilon$、$D_0=\sum_i p_i+\sum_i y_i+\epsilon$，则
\begin{equation}
 \frac{\partial\mathcal L_{\rm Dice}}{\partial p_j}
 =\frac{N-2y_jD_0}{D_0^2},\qquad
 \frac{\partial\mathcal L_{\rm Dice}}{\partial z_j}
 =\frac{N-2y_jD_0}{D_0^2}p_j(1-p_j).
\end{equation}
一个像素的梯度依赖全图预测总量，体现其区域级归一化，而不是各像素独立计分。对无前景图，$N=\epsilon$，当预测面积较大时惩罚梯度可能很小；可结合 BCE 并明确空掩码的评价约定。Dice也不是概率校准目标，因此不应把优化后的$p_i$未经验证就解释为可靠风险概率。

硬掩码下，令$J=\operatorname{IoU}$，由$|A|+|B|=|A\cup B|+|A\cap B|$得到$D=2J/(1+J)$。同一对掩码的两个指标单调对应，但先对图像平均再作非线性变换不等价；软Dice与阈值后二值IoU更不能直接替换。边界和拓扑验收正是补充这些面积指标遗漏的信息。

\subsection{U-Net 的改进与 Mask R-CNN 的实例输出}
U-Net++ 在横向连接路径中增加卷积与中间融合，使解码器逐步组合不同深度的表示，代价是更多计算和存储。DeepLab 系列则用空洞卷积在不继续下采样的情况下扩大上下文：空洞率为二的$3\times3$核覆盖五像素宽的范围，仍只有九个空间采样点；在式\eqref{eq:vision-receptive-field}中将增量改为$d_l(k_l-1)j_{l-1}$，即可计算空洞率$d_l$的感受野。空洞空间金字塔池化（Atrous Spatial Pyramid Pooling，ASPP）并行使用不同空洞率，组合不同范围的上下文。稀疏采样不等于观察了覆盖范围内的全部细节，边界仍要依赖足够密的特征。这些改进是否改善边界，都要用相同数据与评价协议检验。

语义掩码中，同类相邻目标仍可能连在一起。若要分别统计每个缺陷的面积，可以使用 Mask R-CNN\footnote{来源：He 等，\href{https://arxiv.org/abs/1703.06870}{\textit{Mask R-CNN}}，2017。}：它在 Faster R-CNN 中增加一个与分类和框回归并行的掩码分支，用 RoIAlign 按候选框位置在共享特征上插值采样，为每个正候选预测固定网格的掩码，推理时再映射回该框在原图中的位置。它依赖检测结果，漏掉的候选不会被掩码分支补回，较粗的区域网格也会限制细边界。

\section{少样本异常检测}
\label{sec:vision-fewshot-ad}
\subsection{只有少量正常样本时要回答什么}
分类、检测和分割都需要缺陷样本，而新产线或低缺陷率工序常常只有少量正常件，缺陷种类也难以穷举。异常检测只用正常样本建立参照：测试时给整幅图一个异常得分，并为每个位置给出偏离程度，后者称为异常图；得分超过阈值的图送去复检，异常图指出可疑位置。每类产品只提供$k$张正常参考图（如$k=1,2,4$）时，称为少样本（few-shot）异常检测。异常得分衡量的是与参照的差异，不是缺陷概率；异常图高亮的位置也未必都是缺陷。

\subsection{DINOv2 与无标签图像块特征}
要让逐块比较可靠，图像块特征应当对同一部位稳定、对异常变化敏感。DINOv2 是不使用人工标签训练的 ViT\footnote{来源：Oquab 等，\href{https://arxiv.org/abs/2304.07193}{\textit{DINOv2: Learning Robust Visual Features without Supervision}}，2024。}：教师网络取学生网络参数的指数滑动平均，两者看到同一幅图的不同裁剪；学生既要让分类标记的输出与教师一致，也要在被遮住的图像块位置预测教师给出的图像块表示。训练后，不同图像中相同部位的图像块特征彼此相近。本节只把冻结的 DINOv2 当作特征提取器，训练目标留到预训练章展开。它的普通版本有第\ref{sec:vision-register-artifacts}节所说的高范数伪影，官方也提供了带寄存标记的版本；后续的 DINOv3 进一步扩大了模型与数据\footnote{来源：Siméoni 等，\href{https://arxiv.org/abs/2508.10104}{\textit{DINOv3}}，2025。}。

设冻结的骨干把图像$x$变成$N$个图像块特征$f_1(x),\ldots,f_N(x)\in\R^D$。ViT-S/14 的图像块边长为14像素、$D=384$，输入为$448\times448$时有$(448/14)^2=1024$个图像块。

\subsection{以最近邻距离给图像块打分}
AnomalyDINO 直接比较这些图像块特征，不训练任何参数\footnote{来源：Damm 等，\href{https://arxiv.org/abs/2405.14529}{\textit{AnomalyDINO: Boosting Patch-based Few-shot Anomaly Detection with DINOv2}}，2025。}。先把$k$张正常参考图的全部图像块特征放入记忆库
\begin{equation}
 \mathcal M=\{f_j(x^{(r)}):r=1,\ldots,k;\ j=1,\ldots,N\},
\end{equation}
再对测试图的每个图像块求到记忆库的最近余弦距离，并取距离最大的1\%图像块的平均值作为整图得分：
\begin{equation}
 d_j=\min_{m\in\mathcal M}\left(1-\frac{f_j(x)^{\mathsf T}m}{\|f_j(x)\|_2\,\|m\|_2}\right),\qquad
 s(x)=\frac1{|\mathcal T|}\sum_{j\in\mathcal T}d_j,
 \label{eq:vision-fewshot-score}
\end{equation}
其中$\mathcal T$是$d_j$最大的$\lceil0.01N\rceil$个图像块。把$d_1,\ldots,d_N$排回图像块网格，经上采样和平滑便得到异常图。参考图默认加入旋转增强，物体类图像还会先去掉背景，图\ref{fig:vision-fewshot-ad}中把两者标为可选步骤。“正常图像块记忆库加最近邻距离”的思路此前已用于 PatchCore\footnote{来源：Roth 等，\href{https://arxiv.org/abs/2106.08265}{\textit{Towards Total Recall in Industrial Anomaly Detection}}，2022。}；AnomalyDINO 表明，换用 DINOv2 特征后，少到一张参考图也能工作。

\begin{figure}[htbp]
\centering
\begin{tikzpicture}[node distance=0.5cm and 0.5cm]
 \node[idi input, minimum width=2.0cm] (ref) {正常参考图\\$k$ 张};
 \node[idi process, right=of ref, minimum width=1.9cm] (bb1) {冻结的\\DINOv2};
 \node[idi state, right=1.9cm of bb1, minimum width=2.6cm] (mem) {记忆库 $\mathcal M$\\正常图像块特征};
 \node[idi input, below=1.2cm of ref, minimum width=2.0cm] (test) {测试图像\\$x$};
 \node[idi process, right=of test, minimum width=1.9cm] (bb2) {冻结的\\DINOv2};
 \node[idi process, right=1.9cm of bb2, minimum width=2.6cm] (nn) {各图像块的\\最近邻距离 $d_j$};
 \node[idi output, right=of nn, yshift=0.8cm, minimum width=2.3cm] (map) {异常图\\上采样、平滑};
 \node[idi output, right=of nn, yshift=-0.8cm, minimum width=2.3cm] (score) {整图得分\\最大 1\% 的均值};
 \draw[idi arrow] (ref)--(bb1);
 \draw[idi arrow] (bb1)--node[above,font=\scriptsize,align=center]{可选：旋转\\去背景}(mem);
 \draw[idi arrow] (test)--(bb2);
 \draw[idi arrow] (bb2)--node[above,font=\scriptsize]{$N$ 个特征}(nn);
 \draw[idi arrow] (mem)--node[right,font=\scriptsize]{逐块比较}(nn);
 \draw[idi arrow] (nn.east)--(map.west);
 \draw[idi arrow] (nn.east)--(score.west);
\end{tikzpicture}
\caption{基于图像块最近邻的少样本异常检测。上行用少量正常参考图建立记忆库，下行求测试图每个图像块到记忆库的最近距离；骨干只提取特征，参数不更新。}
\label{fig:vision-fewshot-ad}
\end{figure}

\paragraph{手算两个图像块的距离。}
设记忆库只有$m_1=(1,0)$与$m_2=(0.8,0.6)$两个单位向量。测试图像块$a=(0.6,0.8)$与它们的余弦相似度为$0.6$和$0.96$，距离为$0.4$和$0.04$，取最小值得$d_a=0.04$，说明$a$接近某个正常部位。图像块$b=(-0.6,0.8)$的相似度为$-0.6$和$0$，距离为$1.6$和$1$，于是$d_b=1$。若整图有$200$个图像块，$\mathcal T$只含$2$个；最大的两个距离为$1$和$0.6$时，$s=0.8$。只要少数图像块明显偏离，整图得分就会升高，这对小缺陷有利；平均两个最大值又能减弱单个噪声图像块的影响。最近邻取最小值，意味着测试块只要与任一正常块相似就不会被报出，因此记忆库应覆盖正常件允许出现的各种外观。逐块比较的计算量随记忆库规模增长：$448\times448$输入、$k=4$时，每张测试图需计算$1024\times4096\approx4.2\times10^6$个余弦相似度，可写成一次矩阵乘法。

\subsection{适用边界与验证}
最近邻比较只能发现“与参考图不同”的外观。参考图太少时，光照、位置、批次纹理等正常变化也会被报为异常，应补充覆盖这些变化的参考图或先做对位。旋转增强假定旋转后的正常件仍然正常；若方向本身就是检验内容，如零件装反，就不能使用。去背景也可能把物体的一部分误当作背景。尺寸下限取决于图像块大小和输入分辨率：把$2048$像素宽的原图缩放到$448$像素时，每个图像块对应原图约$64\times64$像素，宽$10$像素的裂纹只剩约$2$像素，提高分辨率又会增加计算。缺件、错装等逻辑异常逐块看可能都正常，需要整体比较。部署阈值须在验证集上确定；异常不等于缺陷，最终仍要按第\ref{sec:vision-validation}节的做法用独立数据检验。

\section{质检验证与适用边界}
\subsection{细小缺陷与有限标注}
分类、检测和分割的错误要分别检查。若分类正确而检测漏掉细裂纹，应先看裂纹是否在缩放或下采样时消失，再检查候选分配与筛选。如果框已经正确而掩码边缘断裂，则应检查解码分辨率、像素标签与掩码阈值。换一个更深的骨干不一定解决这些不同来源的问题。

标注不足时，可以使用在其他图像上预训练的骨干。冻结骨干、只训练最终线性分类层，称为线性探测：骨干每图输出$128$维向量时，三类分类头只有$3\times128+3=387$个参数；全量微调还会更新骨干。检测和分割则必须训练各自的空间任务头；若预训练没有保留所需纹理，仅更新分类头可能不足。不用人工类别标签的预训练（如第\ref{sec:vision-fewshot-ad}节的 DINOv2）将在预训练章展开，图文对齐放到多模态章讨论。

数据增强也要服从任务。几何变换必须同步作用于图像、框和掩码。离散类别掩码缩放通常采用最近邻插值。若裁剪删掉全部裂纹，图像级标签也要相应检查。若颜色是缺陷的依据，过强的颜色扰动会破坏标签。增强需要模拟允许出现的变化，不能随意要求模型忽略任务所需信息。

\subsection{用独立数据检查类别、位置和边界}
\label{sec:vision-validation}
设生产线上有四台相机。可以用前三台相机的早期批次训练、后期批次验证，完整留出第四台相机测试设备迁移。来自同一工件的近重复帧应放在同一集合，避免模型在测试时再次看到几乎相同的图。阈值、输入分辨率与后处理规则在验证集确定。固定测试集只用于最终评价。若要区分相机变化与时间变化的影响，还需要分别固定其中一个因素进行测试。

分类除准确率外，还应报告缺陷召回和正常图误报。例如十个真实缺陷中找对八个，另误报两个，精确率为$8/(8+2)=0.8$，召回率为$8/10=0.8$。检测须先完成前述类别与空间匹配，再统计这些数量；分类热力图可以帮助检查模型是否依赖批次标记，却不能替代经过框标注训练与评价的检测结果。分割还要在 IoU、Dice 等面积重叠之外检查连通性和边界距离：细裂纹少量断裂只改变很少像素，却影响长度估计。面积过滤等后处理也可能删除真实小缺陷，应纳入完整预测流程一起评价；标注边缘本身模糊时，可保留多名标注者的分歧，避免把标注不确定性全部归因于模型。

置信度也需要单独验证：在大量被报为$0.8$缺陷概率的图像中，实际缺陷比例接近$80\%$，才支持这一概率解释；新相机与光照变化可能破坏这种对应。无法自动判定的图像可以转人工复核，同时报告自动覆盖率：例如一百张图自动处理八十张、其中两张错误，覆盖率为$80\%$，已处理样本错误率为$2/80=2.5\%$。最后应在目标设备上测量包含预处理与后处理的延迟，并记录相机、批次、标注版本和所用模型，使漏检可以追溯到具体环节。

以下材料按本章顺序排列，可在完成相关算例后阅读或实践。
\section*{辅助学习材料}
\begingroup
\emergencystretch=3em
\begin{itemize}
\item \textbf{Convolutional Neural Networks (CNNs / ConvNets)}（Stanford CS231n，英文笔记）。复习卷积层、输出尺寸与参数共享。\href{https://cs231n.github.io/convolutional-networks/}{在线阅读}。
\item \textbf{ResNet网络讲解}（中文视频）。区分恒等捷径与改变尺寸的投影捷径。\href{https://www.bilibili.com/video/BV1T7411T7wa}{观看讲解}。
\item \textbf{Vision Transformer网络讲解}（中文视频）。看图像如何切成 patch 并投影为 token。\href{https://www.bilibili.com/video/BV1Jh411Y7WQ}{观看讲解}。
\item \textbf{区域卷积神经网络（R-CNN）系列}（《动手学深度学习》，中文教程）。先读 Fast R-CNN 与 RoI pooling，再读引入 RPN 的 Faster R-CNN。\href{https://zh.d2l.ai/chapter\_computer-vision/rcnn.html}{在线阅读}。
\item \textbf{YOLO系列网络讲解(V1\textasciitilde V3)}（中文视频）。比较网格、锚框与多尺度输出。\href{https://www.bilibili.com/video/BV1yi4y1g7ro}{观看讲解}。
\item \textbf{U-Net网络讲解}（中文视频）。沿编码器、解码器与同尺度拼接追踪分辨率。\href{https://www.bilibili.com/video/BV1Vq4y127fB/}{观看讲解}。上述四个模型视频均可从\href{https://github.com/WZMIAOMIAO/deep-learning-for-image-processing}{作者课程目录}查到配套代码。
\item \textbf{TorchVision Object Detection Finetuning Tutorial}（PyTorch 官方，英文实践）。跟随 Mask R-CNN 微调流程，核对框、标签和掩码格式。\href{https://docs.pytorch.org/tutorials/intermediate/torchvision\_tutorial.html}{在线阅读}。
\end{itemize}
\endgroup
\paragraph{本章小结。}
本章从同一幅质检图像的三种输出展开。分类把空间特征汇聚成类别判断，ResNet 用残差路径改善深层卷积网络的优化，ViT 用图像块注意力联系不同区域；寄存标记为内部计算提供额外位置，但分类正确并不保证局部特征与缺陷对齐。

要求位置时，Faster R-CNN 用共享特征、区域提议和区域识别输出检测框，单阶段检测省去逐候选识别，FPN 让不同大小的缺陷有合适的预测层级；要求边界时，U-Net 用解码器与同尺度连接输出像素标签，Mask R-CNN 进一步区分每个目标。缺陷样本稀少时，冻结的 DINOv2 图像块特征可以与少量正常样本逐块比较，给出整图得分与异常图；它衡量的是偏离正常参照的程度，仍需用缺陷样本验证。框匹配、像素损失和最终评价各自回答不同问题，不能用一个训练分数代表全部任务效果。

找到裂纹后，还需要结合设备关系、有效规程和维修记录判断如何处理，后续章节将分别讨论。视觉模型给出的类别、坐标和掩码应保留拍摄条件、尺度和模型版本，便于他人核对。
